\section*{Chapitre \ref{ch:g11:reaction-energy} --- L'énergie des réactions~: combustion et énergies de liaison}

\begin{solution}{exo:g11:reaction-energy:1}
Bois qui \omterm{def:g4:what-a-fire-needs:burning}{brûle}~: \omterm{def:g11:reaction-energy:exothermic}{exothermique}. Glace qui fond~: \omterm{def:g11:reaction-energy:exothermic}{endothermique} (elle
prend de la chaleur au milieu extérieur, bien que ce ne soit pas une
\omterm{def:g7:chemical-reactions:reaction}{réaction chimique}). Poche de froid~: \omterm{def:g11:reaction-energy:exothermic}{endothermique}. Chaufferette~:
\omterm{def:g11:reaction-energy:exothermic}{exothermique}. Fabrication du sucre à la lumière~: \omterm{def:g11:reaction-energy:exothermic}{endothermique}
(l'énergie vient de la lumière).
\end{solution}

\begin{solution}{exo:g11:reaction-energy:2}
$Q = 3.0 \times 890.7 = \qty{2.7e3}{kJ}$, soit environ \qty{2.7}{MJ}.
\end{solution}

\begin{solution}{exo:g11:reaction-energy:3}
Les \omterm{def:g7:chemical-reactions:reactant}{réactifs}. La différence est cédée au milieu extérieur,
principalement sous forme de chaleur.
\end{solution}

\begin{solution}{exo:g11:reaction-energy:4}
$n = 100 / 46.0 = \qty{2.17}{mol}$~; $Q = 2.17 \times 1367.6 =
\qty{2.97e3}{kJ}$, soit environ \qty{3.0}{MJ}.
\end{solution}

\begin{solution}{exo:g11:reaction-energy:5}
C'est l'énergie nécessaire pour rompre une \omterm{def:g10:the-mole:amount}{mole} de ces liaisons, les
\omterm{def:g7:atoms-and-molecules:molecule}{molécules} et les \omterm{def:g7:atoms-and-molecules:atom}{atomes} étant gazeux. Le doublet mis en commun maintient
les deux \omterm{def:g7:atoms-and-molecules:atom}{atomes} liés~: les écarter l'un de l'autre s'oppose à cette
attraction, ce qui demande de l'énergie.
\end{solution}

\begin{solution}{exo:g11:reaction-energy:6}
Rompues~: $2 \times 436 + 498 = \qty{1370}{kJ}$~; formées~:
$4 \times 464 = \qty{1856}{kJ}$~; $E_r \approx 1370 - 1856 =
\qty{-486}{kJ}$~: la réaction est \omterm{def:g11:reaction-energy:exothermic}{exothermique}.
\end{solution}

\begin{solution}{exo:g11:reaction-energy:7}
L'éthanol \ce{CH3-CH2-O-H} possède 5 liaisons \ce{C-H}, 1 \ce{C-C},
1 \ce{C-O} et 1 \ce{O-H}. Rompues~: $5 \times 415 + 345 + 350 + 464 + 3 \times 498 =
\qty{4728}{kJ}$. Formées~: $4 \times 741 + 6 \times 464 = \qty{5748}{kJ}$.
$E_r \approx \qty{-1020}{kJ}$, inférieure d'environ un quart, en valeur
absolue, aux \qty{1367.6}{kJ} mesurés.
\end{solution}

\begin{solution}{exo:g11:reaction-energy:8}
Propane~: $2219.2 / 0.0440 = \qty{5.04e4}{kJ/kg} = \qty{50.4}{MJ/kg}$.
Méthane~: $890.7 / 0.0160 = \qty{55.7}{MJ/kg}$. Les deux valeurs sont
celles du diagramme.
\end{solution}

\begin{solution}{exo:g11:reaction-energy:9}
$Q = 1500 \times 4.18 \times 85 = \qty{5.33e5}{J} = \qty{533}{kJ}$~;
$n = 533 / 890.7 = \qty{0.598}{mol}$, soit $0.598 \times 16.0 =
\qty{9.6}{g}$ de méthane~; si la moitié de l'énergie est perdue,
\qty{19}{g}.
\end{solution}

\begin{solution}{exo:g11:reaction-energy:10}
Dihydrogène > méthane > propane > octane > éthanol.
$48.0 / 142.9 = \qty{0.34}{kg}$ de dihydrogène.
\end{solution}

\begin{solution}{exo:g11:reaction-energy:11}
Rompues~: $946 + 3 \times 436 = \qty{2254}{kJ}$~; formées~:
$6 \times 390 = \qty{2340}{kJ}$~; $E_r \approx \qty{-86}{kJ}$~: la
réaction est légèrement \omterm{def:g11:reaction-energy:exothermic}{exothermique}.
\end{solution}

\begin{solution}{exo:g11:reaction-energy:12}
Eau~: $250 \times 4.18 \times 20 = \qty{2.09e4}{J} = \qty{20.9}{kJ}$.
Éthanol~: $1.20 / 46.0 = \qty{0.0261}{mol}$, qui pourraient libérer
$0.0261 \times 1367.6 = \qty{35.7}{kJ}$. Rendement
$20.9 / 35.7 \approx \qty{59}{\percent}$. Pertes~: chaleur emportée par
l'\omterm{def:g4:air-a-mixture-of-gases:air}{air}, chaleur qui réchauffe la canette et le support, \omterm{def:g8:combustion-and-fuels:complete}{combustion
incomplète} (suie), éthanol qui s'évapore sans \omterm{def:g4:what-a-fire-needs:burning}{brûler}.
\end{solution}

\begin{solution}{exo:g11:reaction-energy:13}
\ce{2C8H18 + 25O2 -> 16CO2 + 18H2O}~: 8 \omterm{def:g10:the-mole:amount}{moles} de dioxyde de carbone,
\qty{352}{g}, pour \qty{5470}{kJ}, donc $352 / 5.470 = \qty{64.4}{g}$
par mégajoule. \ce{C3H8 + 5O2 -> 3CO2 + 4H2O}~: \qty{132}{g} pour
\qty{2219.2}{kJ}, donc \qty{59.5}{g} par mégajoule. Le méthane, avec
\qty{49.4}{g}, en libère le moins.
\end{solution}

\begin{solution}{exo:g11:reaction-energy:14}
La table donne des \omterm{def:g11:reaction-energy:bond-energy}{énergies de liaison} moyennes, alors que les liaisons
\ce{C=O} du dioxyde de carbone sont plus fortes que la moyenne~; les
valeurs mesurées correspondent à de l'eau liquide, dont la condensation
libère plus d'énergie que la réaction entre gaz~; et la méthode traite
chaque liaison comme indépendante de ses voisines. L'estimation reste
utile~: elle donne le bon signe (\omterm{def:g11:reaction-energy:exothermic}{exothermique}) et le bon ordre de
grandeur sans aucune mesure.
\end{solution}

\begin{solution}{exo:g11:reaction-energy:15}
Rompues~: $4 \times 464 = \qty{1856}{kJ}$~; formées~: $2 \times 436 + 498 =
\qty{1370}{kJ}$~; $E_r \approx \qty{+486}{kJ}$~: la réaction est
\omterm{def:g11:reaction-energy:exothermic}{endothermique}, il faut donc fournir de l'énergie. La \omterm{def:g4:what-a-fire-needs:burning}{combustion}
ultérieure du dihydrogène restitue cette énergie~: le dihydrogène stocke
de l'énergie produite ailleurs, il n'en crée pas.
\end{solution}

\begin{solution}{pb:g11:reaction-energy:1}
\textbf{1.} \ce{CH4 + 2O2 -> CO2 + 2H2O}.

\textbf{2.} \ce{C2H6O + 3O2 -> 2CO2 + 3H2O}.

\textbf{3.} \ce{2C8H18 + 25O2 -> 16CO2 + 18H2O}.

\textbf{4.} \ce{2H2 + O2 -> 2H2O}.

\textbf{5.} Le dihydrogène.

\textbf{6.} \qty{16.0}{g/mol}, \qty{46.0}{g/mol}, \qty{114.0}{g/mol},
\qty{2.0}{g/mol}.

\textbf{7.} En \si{kJ} par \si{kg}, puis en \si{MJ/kg}~: méthane
$890.7 / 0.0160$, \qty{55.7}{MJ/kg}~; éthanol $1367.6 / 0.0460$,
\qty{29.7}{MJ/kg}~; octane $5470 / 0.1140$, \qty{48.0}{MJ/kg}~;
dihydrogène $285.8 / 0.0020$, \qty{142.9}{MJ/kg}.

\textbf{8.} Dihydrogène > méthane > octane > éthanol.

\textbf{9.} Le dihydrogène est un gaz très léger~: un kilogramme occupe
un volume énorme à la pression ordinaire, et il faut donc le comprimer
dans de lourds réservoirs à haute pression.

\textbf{10.} Méthane~: 4 \ce{C-H}~; dioxygène~: 2 \ce{O=O} rompues.
Dioxyde de carbone~: 2 \ce{C=O}~; eau~: $2 \times 2 = 4$ \ce{O-H}
formées.

\textbf{11.} $4 \times 415 + 2 \times 498 = \qty{2656}{kJ}$.

\textbf{12.} $2 \times 741 + 4 \times 464 = \qty{3338}{kJ}$.

\textbf{13.} $E_r \approx 2656 - 3338 = \qty{-682}{kJ}$, contre
\qty{-890.7}{kJ} mesurés~: environ \qty{23}{\percent} de moins en valeur
absolue.

\textbf{14.} Les \omterm{def:g11:reaction-energy:bond-energy}{énergies de liaison} sont des moyennes (la liaison
\ce{C=O} du dioxyde de carbone est plus forte que la moyenne)~; l'eau est
liquide dans la mesure, gazeuse dans l'estimation.

\textbf{15.} $1000 / 890.7 = \qty{1.123}{mol}$.

\textbf{16.} $1.123 \times 44.0 = \qty{49.4}{g}$.

\textbf{17.} Éthanol~: $1000 / 1367.6 = \qty{0.731}{mol}$, qui donnent
\qty{1.462}{mol} de dioxyde de carbone, soit \qty{64.3}{g}. Octane~:
$1000 / 5470 = \qty{0.183}{mol}$, qui donnent \qty{1.463}{mol}, soit
\qty{64.4}{g}.

\textbf{18.} Le méthane~: c'est lui qui possède le plus d'\omterm{def:g7:atoms-and-molecules:atom}{atomes}
d'hydrogène par \omterm{def:g7:atoms-and-molecules:atom}{atome} de carbone (4, contre 2.25 pour l'octane), et la
\omterm{def:g4:what-a-fire-needs:burning}{combustion} de l'hydrogène fournit de l'énergie sans dioxyde de carbone.

\textbf{19.} De la façon dont le dihydrogène est produit~: obtenu à
partir de l'eau avec de l'électricité d'origine renouvelable, il ne
libère presque pas de dioxyde de carbone~; produit à partir du méthane,
il en libère à l'usine au lieu du pot d'échappement.

\textbf{20.} Le méthane libère environ \qty{49}{g} de dioxyde de carbone
par mégajoule.
\end{solution}
